\chapter{多尺度物理模型}
\label{ch:physics}

停机、探针和解释至少要看见温度、液相含水和反应进度。中心温度到达某一阈值只是跟踪器的中间量；十六条控制体的工程通量是评分核，连续方程是规格。下面先写出连续形式，再收缩到实际积分的离散核。收缩保留冷点、表面温度、平均含水和致死进度，压力场、蒸汽质量分数和三维流动留在规格里。界面潜热用蒸汽指令和活度的工程式给出，方向与冷凝放热、蒸发吸热一致。

产品侧对应的收缩更直白。说明书不会写出多孔介质的体积平均方程，但会写出“蒸汽冷凝加热”和“热风烘烤”。冷凝是界面上 \(p_v>p_{\mathrm{sat}}\) 时潜热为正，烘烤是蒸汽压低、蒸发项吸热并留下干壳。符号由分压差决定，不由档位名称决定。蒸档把加热夹住、蒸汽可以离开零，于是冷凝项更容易为正；烤档把蒸汽写成零、排汽抬高，于是蒸发项更容易为负。十六条控制体把这块食品沿厚度切开，冷点取最低一层，表面取最外一层。面板上的色带若画成均匀一块，把这两层折成了同一个数。

\section{计算域与全场}

食品区域 \(\Omega_f\subset\R^3\)，腔体气体区域 \(\Omega_c\)，交界面 \(\Gamma=\partial\Omega_f\)。体积分数满足
\[
\varepsilon_s+\varepsilon_l+\varepsilon_g=1,\qquad \varepsilon_s=\varepsilon_s(S)>0.
\]
\(S(\vect{r},t)\) 综合收缩、孔隙闭合与表面结壳。全场状态为
\begin{equation}
\label{eq:state}
\vect{X}(\vect{r},t)=\big(T,\,p_g,\,\varepsilon_l,\,\omega_v,\,\vect{\alpha},\,S\big).
\end{equation}
反应进度 \(\vect{\alpha}=(\alpha_p,\alpha_g,\alpha_b,\alpha_m)\) 依次是蛋白质变性、淀粉糊化、美拉德褐变与微生物致死当量。腔体状态不进入食品场，而作为边界输入 \((T_\infty,Y_\infty,\vect{v}_\infty,\dot m_{\mathrm{st}},P_{\mathrm{heat}})\)。

十六条控制体没有独立的 \(p_g\) 和 \(\omega_v\) 未知量。写出全场是为了标明离散项各自对应连续方程里的哪一项。

\section{能量与三相质量}

忽略热局部非平衡时，固、液、气在一个代表体元内用同一个 \(T\)。该简化在蒸烤的分钟尺度上可接受。能量方程保留四项：
\begin{equation}
\label{eq:energy}
(\rho c)_{\mathrm{eff}}\frac{\partial T}{\partial t}
=\nabla\!\cdot\!\big(k_{\mathrm{eff}}\nabla T\big)
-\sum_{i\in\{l,v,a\}}\nabla\!\cdot\!(h_i\vect{n}_i)
-\dot m_{\mathrm{phase}}\,\Delta h_{\mathrm{vap}}
+q_{\mathrm{reac}}.
\end{equation}
有效容积热容与导热随 \(\varepsilon_l\) 和 \(S\) 变化，表面变干后导热下降，这是外干内生的机制。液相与蒸汽
\begin{align}
\label{eq:liquid}
\frac{\partial(\varepsilon_l\rho_l)}{\partial t}
&=-\nabla\!\cdot\!\vect{n}_l-\dot m_{\mathrm{phase}},\\
\label{eq:vapor}
\frac{\partial(\varepsilon_g\rho_g\omega_v)}{\partial t}
&=-\nabla\!\cdot\!\vect{n}_v+\dot m_{\mathrm{phase}}.
\end{align}
通量由 Darcy 渗流与 Fick 扩散合成。相变保持有限速率：
\begin{equation}
\label{eq:phase}
\dot m_{\mathrm{phase}}
=K_{\mathrm{ev}}\,a_w(\varepsilon_l,T)\,
\big(p_{\mathrm{sat}}(T)-p_v\big)_+
-K_{\mathrm{cd}}\,(p_v-p_{\mathrm{sat}}(T))_+.
\end{equation}
蒸制时腔体蒸汽压高，冷凝项在表面释放潜热；烤制时蒸汽压低，蒸发项吸热并留下干壳。符号由 \(p_v-p_{\mathrm{sat}}\) 决定，不由档位名称决定。

本次核没有压力未知量。界面潜热被收成
\begin{equation}
\label{eq:qlat}
q_{\mathrm{lat}}=18000\,u_{\mathrm{st}}(1-a_w)-9000\,(1-u_{\mathrm{st}})\max(a_w-0.35,0)\,u_{\mathrm{fan}}.
\end{equation}
蒸汽指令直接决定符号，活度只调制幅度。这保留冷凝放热、蒸发吸热的方向，没有保留\eqref{eq:phase}对局部蒸汽压的依赖。

\section{腔体、表面与一维核}

家用蒸烤箱的腔体是一个带玻璃门的封闭盒子。容积常见在 20--30 L，内壁是拉丝不锈钢，底部有带冲孔的烤盘，后壁风机把热风打转，顶部或背部的加热管提供辐射。蒸汽从一侧喷嘴进入，排汽从另一侧百叶离开。水位写在水箱里，门写在铰链微动开关上。数值核没有单独的三维流动场，把这些零件收成两个标量：腔温 \(T_c\) 和腔湿 \(h_c\)。加热指令把热量送进空气，排汽把热量和湿气带走，蒸汽指令把湿气送进来。食品侧看见的是界面上的对流、辐射和潜热，不是喷嘴的几何细节。

腔温与腔湿
\begin{align}
\dot T_c &= 0.02\,u\,(T_h-T_c)-0.01\,e\,(T_c-25),\\
\dot h_c &= 0.015\,s-0.02\,e\,h_c.
\end{align}
换热系数 \(h=12+55f\)。风机从 0.4 升到 1.3，\(h\) 从 34 升到 83.5，单位是瓦每平方米开尔文。水的饱和蒸汽压用 Antoine 式，温度取摄氏度，压力取 bar：
\begin{equation}
\label{eq:antoine}
\log_{10} P^\star = 5.1962 - \frac{1730.63}{T+233.426}.
\end{equation}
腔体侧分压是腔湿乘以腔温下的饱和压，食品侧分压是表面活度乘以表面温度下的饱和压。两侧之差再乘 \(4000\,\mathrm{W}/(\mathrm{m}^2\cdot\mathrm{bar})\)，加到潜热上。辐射没有单独的角系数场，加热指令上保留 \(7000u\) 作为工程份额。

食品沿厚度分成十六条控制体。下标 0 在界面一侧，下标 15 在对称面。时间步 \(\Delta t=2\,\mathrm{s}\)。导热基线 \(k=\alpha\rho c\)。结壳变量 \(\xi\) 在表面含水低于 \(0.45\) 时增长，增长之后表面那一条的导热乘 \(1-0.55\xi\)。含水扩散系数 \(D=2\times 10^{-9}(1+3m)\)。15 mm 的显式稳定上限约 \(2.20\,\mathrm{s}\)，选用的 \(2\,\mathrm{s}\) 低于该上限。离散方程的完整形式见附录\ref{app:discrete}。

十六条是工程选择，不是物理上的唯一离散。一块 30 mm 的肉糜沿厚度被切成约 1.9 mm 一层，冷点取这十六条里的最低温度。八格会把最低温度所在的那一格加厚，读数偏高，30 mm 对偶的停机偏差会翻号，见第\ref{ch:grid}章。厨房里对应的几何更不规则：肉糜有脂肪夹层，面团有气孔，土豆有表皮。核用一层或两层导热乘 0.45 来代表夹心，用三种食品名来改活度尺度。停止判据读的仍是冷点进度，三种名称的中位数几乎重合。

\begin{figure}[htbp]
\centering
\includegraphics[width=0.86\textwidth]{P_profiles.png}
\caption{种子 7、肉糜 30 mm、组合档、对偶策略：若干时刻的厚度剖面。表面先升，冷点滞后；后期加热下降后剖面重新拉平。}
\label{fig:profiles}
\end{figure}

图\ref{fig:profiles}是同一格式重放得到的内部场。0 min 接近均匀初温；2.5 min 表面已经超过 \(80^\circ\mathrm{C}\)，冷点仍接近初值；10 min 梯度最陡；28 min 接近停机，冷点越过 \(70^\circ\mathrm{C}\)；55 min 加热已经降到保持，剖面接近等温。CSV 只保存表面和冷点最小值，十六条剖面由同一离散格式重放，与日志的最大差低于 \(5\times 10^{-6}\,\mathrm{K}\)。

图\ref{fig:xt}把同一剖面沿时间展开。探针窗口约在 2.5--3.2 min，近表面留下短暂的高温条带，尚未传到冷点。横向没有第二维梯度，这是一维核的几何。图\ref{fig:field3dP}用同一剖面按深度比的 1.15 次幂着色，供演示页面使用；指数是显示约定。含水剖面见图\ref{fig:moistP}：平均含水从 0.72 缓降，结壳门槛 0.45 在 55 min 内未触及，路径 C 因此在这条轨迹上接近零。

\begin{figure}[htbp]
\centering
\includegraphics[width=0.88\textwidth]{P_field_xt.png}
\caption{同一条轨迹的厚度--时间温度场。探针窗口在近表面留下短暂高温条带，尚未传到冷点。}
\label{fig:xt}
\end{figure}

\begin{figure}[htbp]
\centering
\includegraphics[width=0.82\textwidth]{P_field_3d.png}
\caption{十六条剖面的插值色带。着色指数 1.15 是显示约定，横向没有第二维梯度。}
\label{fig:field3dP}
\end{figure}

\begin{figure}[htbp]
\centering
\includegraphics[width=0.82\textwidth]{X_moist.png}
\caption{种子 7 平均含水与表面含水。结壳门槛 0.45 未触及。}
\label{fig:moistP}
\end{figure}

图\ref{fig:paw}给出核内活度对含水。S1 的三档初含水 \(0.64/0.72/0.80\) 都落在饱和段 \(a_w\approx 0.99\)，竖线标出这三档。界面潜热对含水的导数在这一段接近零，这是第\ref{ch:ablate}章含水零效应的函数形状原因。活度是含水的单调饱和函数，与公开 GAB 等温线同方向：含水升高则活度升高并趋近 1。

\begin{figure}[htbp]
\centering
\includegraphics[width=0.82\textwidth]{P_aw.png}
\caption{核内活度。竖线是 S1 三档初含水，三者都在饱和段。}
\label{fig:paw}
\end{figure}
\FloatBarrier

\section{反应进度}

速率形式为
\begin{equation}
\label{eq:kinetics}
\frac{\partial\alpha_i}{\partial t}
=A_i\exp\!\Big(-\frac{E_{a,i}}{RT}\Big)
(1-\alpha_i)^{n_i}\,g_i(\varepsilon_l,S).
\end{equation}
\(g_p\) 在自由水不足时下降，\(g_b\) 在低水分、高温时上升，使表面褐变和内部变性可以分离。热力进度不另造机理：
\begin{equation}
\label{eq:F}
\frac{\mathrm{d}F}{\mathrm{d}t}=10^{(T_{\mathrm{cold}}-T_{\mathrm{ref}})/z},\qquad
\alpha_m=1-\exp(-F/F_{\mathrm{req}}).
\end{equation}
\(T_{\mathrm{cold}}(t)=\operatorname{ess\,inf}_{\vect{r}\in\Omega_f}T(\vect{r},t)\)。核取十六条的最低温度。\(\alpha_m=1-\mathrm{e}^{-F/2}=0.85\) 给出 \(F=-2\ln 0.15=3.794\,\mathrm{min}\)。冷点若保持恒温 \(T\)，每分钟累积 \(10^{(T-70)/10}\) 分钟的等效时间。表\ref{tab:constT}只使用这条映射，不使用任何一条轨迹。

设计线是策略比较的零点，不是认证阈值。面板上若把停机写成“中心到达 \(70^\circ\mathrm{C}\)”，给出的是跟踪器的中间量。同一条 \(70^\circ\mathrm{C}\) 恒温，累积 \(F\) 的速度是每分钟 1，走到 0.85 需要 3.794 min；若冷点只到 \(60^\circ\mathrm{C}\)，每分钟只累积 0.1，走完同一条线要将近 38 min。厨房里“看起来已经熟了”往往读的是表面褐变和蒸汽量，冷点进度可以仍停在 0.1 以下。50 mm 在 55 min 终点的 \(\alpha_m=0.068\) 就是这种情形：表面已经热，内部的等效时间几乎还没有开始积。微生物初始数量不进入状态，食品安全限值需要独立的型式试验。交叉表比较的是四条策略相对这条线的位置。

家用说明书很少写出 \(F\) 值。更常见的是“中心 \(70^\circ\mathrm{C}\)”“表面金黄”“锁住水分”三句并排。三句对应三个不同的量。中心温度是跟踪器的中间量，表面金黄是褐变代理 \(\alpha_b\)，锁住水分是表面含水。停机读的是 \(\alpha_m\)，由冷点温度历史积分而来。种子 7 在 1680 s 冷点已经到 \(70.21^\circ\mathrm{C}\)，进度只有 0.558：温度阈值跨过了，致死当量还没有走到 0.85。再过约三分钟，1866 s 进度才到 0.995。面板若在 1680 s 响铃，给出的是温度阈值，不是设计线。恒温表把这件事写成数字：\(70^\circ\mathrm{C}\) 还要再保持 3.794 min，\(65^\circ\mathrm{C}\) 要 12 min，\(60^\circ\mathrm{C}\) 要将近 38 min。厚件开环在菜单终点往往还停在表的左端。

\begin{table}[htbp]
\centering
\caption{恒温冷点到达 \(\alpha_m=0.85\) 的时间}
\label{tab:constT}
\begin{tabular}{rrrr}
\toprule
\(T/^\circ\mathrm{C}\) & 每分钟的 \(F\) & 所需时间 / min & 所需时间 / s \\
\midrule
55 & 0.0316 & 119.98 & 7199 \\
60 & 0.100 & 37.94 & 2277 \\
65 & 0.316 & 12.00 & 720 \\
70 & 1 & 3.794 & 228 \\
75 & 3.162 & 1.200 & 72 \\
80 & 10 & 0.379 & 23 \\
\bottomrule
\end{tabular}
\end{table}

\(60^\circ\mathrm{C}\) 的 38 min 还放得进 55 min 的时域；\(55^\circ\mathrm{C}\) 的 120 min 放不进。厚件开环在菜单终点冷点往往还在 \(70^\circ\mathrm{C}\) 以下，进度误差贴着 \(-0.85\)，与这张表的左端一致。薄件若把冷点推过 \(80^\circ\mathrm{C}\) 并保持几十秒，进度就会冲到上沿 \(+0.15\)。表不区分这两种情形是哪一种策略造成的。恒温表不是种子 7 的冷点历史，也不能把 0.85 读成法规限值。

结构演化
\begin{equation}
\label{eq:structure}
\frac{\partial S}{\partial t}
=\kappa_s(1-S)\,\langle\alpha_p-\alpha_s^\star\rangle_+
+\kappa_c\,\langle M_{c}- \varepsilon_l\rangle_+.
\end{equation}
核里结壳只在表面含水低于 0.45 时以 \(1.5\times 10^{-4}\,\mathrm{s}^{-1}\) 增长。S1 的初始含水不低于 0.64，表面失水率很低，结壳常停在 0。规格中的路径 C 因此在多数样本上没有闭合。

三个原型共用同一套导热和进度公式。差别只在活度尺度：肉糜 \((0.10,0.70)\)，土豆 \((0.055,0.80)\)，面团 \((0.08,0.75)\)。交叉表里三个食材的中位误差几乎相同，与“原型只改了一个弱的活度尺度”一致。这不能用来声称蛋白与淀粉的反应差异已经被检验。面团膨胀不在状态里。

Kossovich 数
\[
\mathrm{Ko}
=\frac{\Delta h_{\mathrm{vap}}\,\varepsilon_{l0}\,\rho_l}{(\rho c)_{\mathrm{eff}}\,\Delta T}
\]
停在两个未定因子上。核文件里写着潜热常数 \(\Delta h_{\mathrm{vap}}=2.3\times 10^{6}\,\mathrm{J/kg}\)。步进不读取这个常数：界面潜热走的是 \(18000\,u_{\mathrm{st}}\) 与 \(9000\,u_v\) 两条工程通量。体积热容 \((\rho c)_{\mathrm{eff}}=3.4\times 10^{6}\,\mathrm{J}/(\mathrm{m}^3\mathrm{K})\) 进入了温升。两者之比是 \(0.6765\,\mathrm{K\cdot m}^3/\mathrm{kg}\)，于是 \(\mathrm{Ko}=0.6765\,\varepsilon_{l0}\rho_l/\Delta T\)。核里的含水是夹在 \([0.02,0.95]\) 的分数，不是液相体积分数乘密度。因此 Ko 留在这两个因子上，不化成一个纯数。

\section{两节点工作方程}

探针分析使用表面节点 1 与内部节点 2：
\begin{align}
\label{eq:twonode}
C_1\dot T_1
&=\frac{T_\infty-T_1}{R_{\mathrm{ext}}(v)}+\frac{T_2-T_1}{R_{\mathrm{int}}}
+q_{\mathrm{lat}}(\dot m_\Gamma),\\
C_2\dot T_2
&=\frac{T_1-T_2}{R_{\mathrm{int}}},\\
\dot m
&=-k_m\big(a_w(m)\,p_{\mathrm{sat}}(T_1)-p_\infty\big).
\end{align}
\(R_{\mathrm{ext}}\) 由风速与是否冷凝切换，\(R_{\mathrm{int}}\) 随结壳增大。无相变且 \(\dot m=0\) 时，平衡迫使 \(T_1=T_2=T_\infty\)，三个参数从代数方程中消失。把含水冻结，相变通量为零，对 \((T_1,T_2)\) 线性化，
\[
\vect{A}=\begin{pmatrix}-a-b & b\\ c & -c\end{pmatrix},\qquad
\vect{C}=\begin{pmatrix}1 & 0\end{pmatrix},
\]
其中 \(a=1/(C_1 R_{\mathrm{ext}})\)，\(b=1/(C_1 R_{\mathrm{int}})\)，\(c=1/(C_2 R_{\mathrm{int}})\)。可观测矩阵行列式等于 \(b\neq 0\)，两个温度状态从表面看是可观测的。平衡点上对 \(R_{\mathrm{int}}\) 的灵敏度为零。状态可观测给出温度节点可以从表面重建；内部热阻的辨识还需要离开平衡的输入，让灵敏度列不再是零。

说明常数 \(C_1=C_2=2000\,\mathrm{J/K}\)，\(R_{\mathrm{ext}}=0.05\,\mathrm{K/W}\)，\(R_{\mathrm{int}}=0.15\,\mathrm{K/W}\) 时，特征根对应时间常数 \(69.72\,\mathrm{s}\) 与 \(430.3\,\mathrm{s}\)。40 s 的探针窗短于这两个时间常数，表面记录在代数上满秩，数值上仍然弱：窗口结束时冷点几乎还停在初值，信息主要写在表面节点的瞬态上。

界面上还有一项工程辐射份额 \(7000u\)，加在对流与潜热之后。它没有角系数场，也没有发射率随含水变化。种子 7 的表面峰值约 \(150^\circ\mathrm{C}\) 来自这一项与对流的叠加。Antoine 蒸汽压差另乘 \(4000\,\mathrm{W}/(\mathrm{m}^2\cdot\mathrm{bar})\)，符号随腔温与表面温度的相对位置翻转，第\ref{ch:result}章用 216 对积分给出厚度上的符号变化。

评分循环始终积分十六条控制体。压力场、蒸汽质量分数和有限速率相变源项写在连续形式里，用来标明离散项各自对应哪一项；核里看见的是温度、含水、结壳代理和进度。两节点工作方程只用于探针分析：它给出平衡点上内部热阻灵敏度为零、两段脉冲提供第二个灵敏度方向这两条代数事实。第\ref{ch:est}章把跟踪器接到表面读数上，第\ref{ch:ati}章把 40 s 窗口写成已经运行的时间表。

交叉表、网格加密和种子 7 轨迹都读这一组工程常数：扩散率 \(\alpha=2.0\times 10^{-7}\,\mathrm{m}^2/\mathrm{s}\)，容积热容 \(\rho c=3.4\times 10^{6}\,\mathrm{J}/(\mathrm{m}^3\mathrm{K})\)，基线导热 \(k=0.68\,\mathrm{W}/(\mathrm{m\cdot K})\)，时间步 \(2\,\mathrm{s}\)，时域 55 min，控制体数 16；换热 \(h=12+55u_{\mathrm{fan}}\)，界面潜热按\eqref{eq:qlat}，辐射份额 \(7000u\)，Antoine 压差再乘 \(4000\,\mathrm{W}/(\mathrm{m}^2\cdot\mathrm{bar})\)；脂肪夹层把指定控制体的导热乘 0.45。蒸档把加热夹在 0.25，冷凝潜热从外向内送；烤档关掉蒸汽、抬高排汽，管壁辐射和热风把表面推高，内部仍靠导热；组合档放开四个量，探针窗先拉陡表面梯度，窗外回到品质加热。这些数字一旦改写，已经写出的日志不再对应。

体积热容进入温升，界面潜热走的是 \(18000\,u_{\mathrm{st}}\) 与 \(9000\,u_v\) 两条工程通量。核文件里写着 \(\Delta h_{\mathrm{vap}}=2.3\times 10^{6}\,\mathrm{J/kg}\)，步进并不读取这个常数。产品侧对应得很具体：蒸汽发生器的额定克数改的是阀侧流量，食品每一层看见的仍是工程式给出的界面热流。说明书把“蒸汽发生量”写在参数表第一页，发生量改腔湿和界面分压，改不了半厚平方。Kossovich 数因此留在两个未定因子上。要用 Ko 去对标另一台机器，必须先声明液相体积分数和温升取哪一段。本次交叉表没有这一行。

烤盘冲孔、搁架高度和门缝漏汽，在拆机时分属三个零件。核把它们收成换热系数 \(h=12+55u_{\mathrm{fan}}\) 和排汽项。风机从 0.4 升到 1.3，\(h\) 从 34 升到 83.5。探针段故意把风机拧到 1.3，为的是改 \(R_{\mathrm{ext}}\)，不是为了把肉吹熟。操作者若把这一段听成“预热风机”，会把信息动作读成品质动作。日志里模式整数在 150--190 s 是 1，窗外回到 0。同一台风机，两种函数。家用腔体本身也不是规则长方体：搁架挡住一截气流，门缝漏进一股冷空气，冲孔让底部换热和侧面不一样。十六条控制体把这些不规则性压成厚度方向上的一串温度。冷点取最低一层，表面取最外一层。面板上的色带若画成整块均匀肉色，就把这两层折成了同一个数。

辐射份额 \(7000u\) 没有角系数场。种子 7 的表面峰值约 \(150^\circ\mathrm{C}\) 来自这一项与对流的叠加。说明书里的“上下管独立控温”改的是这一份额的上下分配，核把它收成一个工程常数。Antoine 压差再乘 \(4000\,\mathrm{W}/(\mathrm{m}^2\cdot\mathrm{bar})\)，符号随腔温与表面温度的相对位置翻转。蒸档更容易让压差为正，烤档更容易为负。组合档在探针窗内拉满蒸汽，窗后回到 0.15，符号可以在同一次过程里翻转一次。交叉表用 216 对积分给出厚度上的符号变化：15 mm 上移，50 mm 下移。产品文案若把蒸汽一律写成“锁水加热”，就把这一翻转盖住了。

说明书把传热收成三句短话：蒸汽冷凝加热、热风循环、上下管独立。三句各自对应核里的一项，对不上的那一截才需要单独说。冷凝加热对应界面潜热\eqref{eq:qlat} 的正项，阀侧蒸汽指令拉开之后，\(u_{\mathrm{st}}\) 把潜热送进表面格子；热风循环对应 \(h=12+55u_{\mathrm{fan}}\)，风机从 0.4 拧到 1.3，换热从 34 升到 83.5；上下管独立对应辐射份额 \(7000u\)，核里没有角系数场，上下分配被收成同一个工程常数。操作者在面板上拧的是档位和温度，机器改的是这三项的相对大小。冷点进度仍由半厚平方和这一次过程里看见的热阻决定。三句短话闭合的是边界能力，不是停机几何。

上下管在拆机时是两根零件。上管更靠近食品上表面，下管隔着烤盘。核把二者收成同一个辐射份额 \(7000u\)。面板上的“上下管独立”改的是这一份额的分配，积分器没有把分配写成第二个状态。操作者若只开上管，表面先上色、底部仍湿；只开下管，烤盘先热、上表面靠对流跟上。交叉表比较的是四个执行量的开度，不是上下管各自的功率比。把说明书里的独立控温读成两条加热指令，会多出一个核里没有的旋钮。

热风循环那一句也容易被读成“把肉吹熟”。风机改的是边界层厚度。探针段把它拧到 1.3，\(h\) 从 34 升到 83.5，Biot 数从 1.50 升到 3.68，表面梯度更陡，内部更跟不上。买到的是外部热阻和内部热阻的可分性，不是把冷点吹热。种子 7 在这 40 s 里冷点只从 \(8.01^\circ\mathrm{C}\) 升到 \(8.05^\circ\mathrm{C}\)。厨房里风机噪声升高，操作者听见的是叶片加速；十六条里看见的是近表面那一层被拉陡。冷点进度仍按半厚平方走。

家用腔体的零件把这三项钉在看得见的位置。蒸汽从左侧或后壁喷嘴进入，冷凝水沿门玻璃内侧往下走，托盘接住滴水。热风从后壁护网打出来，搁架把一截气流挡住，冲孔烤盘让底部换热和侧面不一样。上下管分别在顶板和底板附近，管壁发红时表面先亮。十六条控制体沿厚度切开这块食品，下标 0 贴着界面，下标 15 贴着对称面。喷嘴位置、搁架遮挡和门缝漏汽不进入这十六条。二维截面把横向那一截重新打开，最低温度被绝热边界钉在最内侧，偏移 \(13.75\,\mathrm{mm}\)。评分循环仍读一维最低一层。产品叙述若把“热风循环更均匀”写成冷点已经居中，截面给出的偏移会把这句话拆开。

拆机时还能看见另外几件说明书常画、积分器不单独成状态的零件。门灯装在铰链侧或顶板前方，开门即亮，用来看上色；核没有光学通道，照明不进入十六条。除垢口在蒸汽发生器一侧，钙化会把阀侧发生量的长期均值拉低，核把发生量收成 \(u_{\mathrm{st}}\) 的瞬时开度。底部冷凝托盘要定期抽出，抽的时候门会开几毫米，微动落下。五档搁架的层距大约两三厘米，上层更靠近加热管，下层更靠近蒸汽喷嘴的冷凝区。核把层高收进同一条对流 \(h=12+55u_{\mathrm{fan}}\)。操作者换一层搁架，改的是辐射份额和对流遮挡的相对大小，半厚平方仍按扩散走。
